\epanalipsi{Γ2}

\tmima{Λεξικό του κεφαλαίου}

\renewcommand{\arraystretch}{1.5}
\begin{longtable}{@{}>{\raggedright\arraybackslash}p{0.28\linewidth} >{\raggedright\arraybackslash}p{0.62\linewidth}@{}}
\toprule
\textbf{Λέξη ή σύμβολο} & \textbf{Με μία πρόταση}\\
\midrule
\endhead
μέση τιμή & Άθροισμα των δεδομένων διά το πλήθος τους.\\
διάμεσος & Η μεσαία τιμή στη σειρά (ή ο μέσος όρος των δύο μεσαίων).\\
μέση τιμή από πίνακα & Άθροισμα των γινομένων τιμή $\times$ συχνότητα, διά το σύνολο των παρατηρήσεων.\\
κέντρο κλάσης & Το μέσο των δύο άκρων της κλάσης.\\
εύρος $R$ & Μεγαλύτερη τιμή μείον μικρότερη τιμή.\\
κορυφή & Η τιμή με τη μεγαλύτερη συχνότητα (μπορεί να υπάρχουν πολλές).\\
απόμακρη τιμή & Τιμή που διαφέρει σημαντικά από τις υπόλοιπες.\\
μεταβλητότητα & Πόσο απλώνονται και διαφέρουν οι τιμές μεταξύ τους.\\
\bottomrule
\end{longtable}

\tmima{Ασκήσεις επανάληψης}

\ex Βρες τη μέση τιμή των $7, 9, 12, 4, 8$.

\ex Βρες τη διάμεσο των $11, 3, 8, 6, 15$ και των $7, 2, 9, 4$.

\ex Δώδεκα παρατηρήσεις έχουν μέση τιμή $6{,}5$. Πόσο είναι το άθροισμά τους;

\ex Η Ζωή έγραψε $13, 16, 14, 15$ στα τέσσερα πρώτα διαγωνίσματα. Τι βαθμό χρειάζεται στο πέμπτο ώστε ο μέσος όρος των πέντε να είναι $15$;

\ex Ο πίνακας δείχνει πόσες ώρες γυμναστικής κάνουν την εβδομάδα $20$ μαθητές. Βρες τη μέση τιμή.\par
\begin{center}
\renewcommand{\arraystretch}{1.25}
\begin{tabular}{lcccc}
\toprule
Ώρες & $1$ & $2$ & $3$ & $4$\\
Μαθητές & $5$ & $8$ & $4$ & $3$\\
\bottomrule
\end{tabular}
\end{center}

\ex Στον πίνακα φαίνονται τα λεπτά που περνούν $20$ παιδιά στο κινητό κάθε μέρα. Βρες τη μέση τιμή με τα κέντρα των κλάσεων.\par
\begin{center}
\renewcommand{\arraystretch}{1.25}
\begin{tabular}{lccc}
\toprule
Λεπτά & $0$--$20$ & $20$--$40$ & $40$--$60$\\
Παιδιά & $4$ & $10$ & $6$\\
\bottomrule
\end{tabular}
\end{center}

\ex Οι ηλικίες επτά αθλητών είναι $14, 18, 14, 22, 17, 14, 19$. Βρες το εύρος και την κορυφή.

\ex Οι βαθμοί έξι μαθητών είναι $15, 16, 14, 15, 2, 16$. Ποια είναι η απόμακρη τιμή; Βρες τη μέση τιμή με και χωρίς αυτή.

\ex Δύο παίκτες πέτυχαν πόντους σε πέντε αγώνες. Ο πρώτος: $2, 10, 6, 14, 8$. Ο δεύτερος: $7, 8, 9, 8, 8$. Βρες τη μέση τιμή και το εύρος του καθενός και σχολίασε.

\ex Το ιστόγραμμα δείχνει πόσες ώρες την εβδομάδα διαβάζουν $20$ μαθητές.
\begin{center}
\binhist{1,4,7,5,3}{0,2,4,6,8,10}{8}{ώρες την εβδομάδα}
\end{center}
Βρες το εύρος, την κλάση-κορυφή και τη μέση τιμή.

\ex Δέκα αριθμοί έχουν μέση τιμή $8$. Προσθέτουμε τον $19$. Ποια είναι η νέα μέση τιμή;

\ex Στο τμήμα Α φοιτούν $5$ μαθητές με μέσο όρο $12$ και στο τμήμα Β $15$ μαθητές με μέσο όρο $16$. Ποιος είναι ο μέσος όρος και των δύο μαζί;

\ex Βρες έξι αριθμούς με διάμεσο $5{,}5$, εύρος $8$ και μέση τιμή $6$.

\ex Οι θερμοκρασίες έξι ημερών ($^\circ$C) είναι $-1, 2, -3, 4, 0, 2$. Βρες μέση τιμή, διάμεσο, εύρος και κορυφή.

\ex Σωστό ή λάθος; α) Η διάμεσος επτά αριθμών είναι ο μέσος όρος των τριών μεσαίων. β) Τέσσερις αριθμοί με άθροισμα $20$ έχουν μέση τιμή $5$. γ) Τα δεδομένα $3,3,5,5,7$ έχουν δύο κορυφές. δ) Οι απόμακρες τιμές δεν επηρεάζουν τη μέση τιμή. ε) Το εύρος των $-5, 0, 4$ είναι $1$.

\ex Σωστό ή λάθος; α) Δύο σύνολα με την ίδια μέση τιμή έχουν πάντα το ίδιο εύρος. β) Όσο μεγαλύτερο το εύρος, τόσο μεγαλύτερη η μεταβλητότητα. γ) Σε ομαδοποιημένα δεδομένα το εύρος είναι το μεγαλύτερο άκρο της τελευταίας κλάσης μείον το μικρότερο άκρο της πρώτης. δ) Αν προσθέσουμε $5$ σε όλες τις τιμές, το εύρος αυξάνεται κατά $5$. ε) Η κορυφή μπορεί να είναι ένα χρώμα.

\tmima{Χάρτης επανάληψης}

\begin{center}
\renewcommand{\arraystretch}{1.3}
\begin{tabular}{ll}
\toprule
\textbf{Ασκήσεις} & \textbf{Παράγραφος}\\
\midrule
1, 2, 3, 4, 5, 6, 11, 12 & Γ2.1 Μέση τιμή και διάμεσος\\
7, 8, 9, 10, 13, 14 & Γ2.2 Εύρος, κορυφές και απόμακρες τιμές\\
15, 16 & Γ2.1 και Γ2.2\\
\bottomrule
\end{tabular}
\end{center}
